\documentclass[12pt,a4paper]{ctexart}

% ===== 页面设置 =====
\usepackage[top=2.5cm,bottom=2.5cm,left=2.5cm,right=2.5cm,headheight=15pt]{geometry}
\usepackage{amsmath,amssymb,amsthm}
\usepackage{enumitem}
\usepackage{xcolor}
\usepackage{tikz}
\usepackage{hyperref}
\usepackage{fancyhdr}

% ===== 页眉页脚 =====
\pagestyle{fancy}
\fancyhf{}
\fancyhead[L]{\textbf{高等数学（第七版·下册）}}
\fancyhead[R]{\textbf{习题11-1 第4题}}
\fancyfoot[C]{\thepage}
\renewcommand{\headrulewidth}{0.8pt}

% ===== 自定义环境 =====
\newtheoremstyle{myremark}
  {0.5\topsep}{0.5\topsep}
  {\normalfont}{}
  {\bfseries\color{blue!60!black}}{：}{0.5em}{}
\theoremstyle{myremark}
\newtheorem*{remark}{解}

% 方框答案
\newcommand{\answerbox}[1]{%
  \begin{center}
  \fbox{\color{red}\large\bfseries #1}
  \end{center}
}

% ===== 文档信息 =====
\title{\Huge\bfseries 习题11-1~~第4题}
\author{高等数学（同济大学数学系~编·第七版·下册）}
\date{\today}

\begin{document}

\maketitle
\thispagestyle{fancy}

% ==================== 原题 ====================
\section*{题目}
\fbox{%
  \parbox{\dimexpr\textwidth-2\fboxsep-2\fboxrule\relax}{%
    求半径为 $a$、中心角为 $2\varphi$ 的均匀圆弧（线密度 $\mu=1$）的质心。%
  }%
}
\bigskip

% ==================== 分析 ====================
\section*{分析}

均匀物体的质心（形心）坐标由以下公式给出：
\[
\bar{x} = \frac{1}{M} \int_L x \, ds,\qquad
\bar{y} = \frac{1}{M} \int_L y \, ds,
\]
其中 $M = \int_L \mu \, ds$ 为物体的总质量。本题中线密度 $\mu=1$，故
\[
M = \int_L 1 \, ds = \text{圆弧的弧长} = a \cdot 2\varphi = 2a\varphi.
\]

\textbf{坐标系的选取是简化计算的关键：}将圆弧对称地放置在 $x$ 轴两侧（如\,
图~\ref{fig:arc}），利用对称性可直接得出 $\bar{y}=0$。

% ==================== 示意图 ====================
\begin{figure}[htbp]
\centering
\begin{tikzpicture}[scale=2.5]
  % 坐标轴
  \draw[->,gray] (-1.5,0) -- (1.8,0) node[right] {$x$};
  \draw[->,gray] (0,-1.3) -- (0,1.3) node[above] {$y$};
  % 圆弧
  \draw[thick,blue] (1,0) arc (0:60:1);
  \draw[thick,blue] (1,0) arc (0:-60:1);
  % 半径虚线
  \draw[dashed,gray] (0,0) -- (60:1) node[midway,above right] {$a$};
  \draw[dashed,gray] (0,0) -- (-60:1);
  % 角度标注
  \draw[->] (0.35,0) arc (0:60:0.35) node[midway,right] {$\varphi$};
  \draw[->] (0.35,0) arc (0:-60:0.35) node[midway,right] {$-\varphi$};
  % 质心
  \fill[red] (0.827,0) circle (0.03) node[below=3pt] {$\bigl(\bar{x},0\bigr)$};
  % 标注
  \node at (0,0) [below left] {$O$};
  \node[blue] at (60:1.15) {$L$};
\end{tikzpicture}
\caption{圆弧 $L$ 对称地置于 $x$ 轴两侧，圆心角为 $2\varphi$}
\label{fig:arc}
\end{figure}

% ==================== 解答 ====================
\section*{解答}

\begin{remark}
  ~
\end{remark}

\medskip

\noindent\textbf{总质量 $M$（即弧长）：}
\[
M = \int_L 1 \cdot ds = \int_{-\varphi}^{\varphi} a \, d\theta
    = a \cdot 2\varphi = 2a\varphi.
\]

\medskip

\noindent\textbf{参数方程：} 取圆心为原点，使圆弧关于 $x$ 轴对称：
\[
x = a\cos\theta,\quad y = a\sin\theta,\qquad -\varphi \leq \theta \leq \varphi.
\]

弧长微元：
\[
ds = \sqrt{\Bigl(\frac{dx}{d\theta}\Bigr)^2 + \Bigl(\frac{dy}{d\theta}\Bigr)^2}\, d\theta
    = \sqrt{(-a\sin\theta)^2 + (a\cos\theta)^2}\, d\theta
    = a\, d\theta.
\]

\medskip

\noindent\textbf{求 $\bar{x}$：}
\[
\int_L x \, ds = \int_{-\varphi}^{\varphi} a\cos\theta \cdot a\, d\theta
                = a^2 \int_{-\varphi}^{\varphi} \cos\theta\, d\theta
                = a^2 \Bigl[\sin\theta\Bigr]_{-\varphi}^{\varphi}
                = a^2 (\sin\varphi - \sin(-\varphi))
                = a^2 \cdot 2\sin\varphi
                = 2a^2 \sin\varphi.
\]
\[
\bar{x} = \frac{1}{M} \int_L x \, ds
        = \frac{2a^2 \sin\varphi}{2a\varphi}
        = \frac{a \sin\varphi}{\varphi}.
\]

\medskip

\noindent\textbf{求 $\bar{y}$（由对称性直接得出）：}

圆弧关于 $x$ 轴对称，被积函数 $y$ 是奇函数（关于 $\theta$），积分区间 $[-\varphi,\varphi]$ 对称：
\[
\int_L y \, ds = \int_{-\varphi}^{\varphi} a\sin\theta \cdot a\, d\theta
                = a^2 \Bigl[-\cos\theta\Bigr]_{-\varphi}^{\varphi}
                = a^2 (-\cos\varphi + \cos(-\varphi))
                = a^2 (-\cos\varphi + \cos\varphi)
                = 0.
\]
\[
\bar{y} = \frac{1}{M} \int_L y \, ds = 0.
\]

\bigskip

\answerbox{$\displaystyle \bar{x} = \frac{a\sin\varphi}{\varphi},\qquad \bar{y} = 0$}

\bigskip

% ==================== 讨论 ====================
\section*{特例验证}

\begin{enumerate}[leftmargin=*]

  \item \textbf{半圆（$\varphi = \dfrac{\pi}{2}$）：}
  \[
  \bar{x} = \frac{a \sin\frac{\pi}{2}}{\frac{\pi}{2}} = \frac{a \cdot 1}{\pi/2} = \frac{2a}{\pi}.
  \]
  这正是半圆形均匀细丝质心的标准结果。

  \item \textbf{整圆（$\varphi = \pi$）：}
  \[
  \bar{x} = \frac{a \sin\pi}{\pi} = 0.
  \]
  完整的均匀圆环质心在圆心，符合对称性预期。

  \item \textbf{当 $\varphi \to 0$（弧长极短）：}
  由 $\displaystyle \lim_{\varphi \to 0} \frac{\sin\varphi}{\varphi} = 1$，得
  \[
  \bar{x} \to a.
  \]
  此时圆弧退化为圆周上一点，质心即在该点位置，符合物理直觉。

\end{enumerate}

\section*{实用意义}

本题的结果 $\bar{x} = \dfrac{a\sin\varphi}{\varphi}$ 是工程力学中的一个重要公式，
在计算曲线形构件的质心、转动惯量等问题中经常用到。它与本节的例2（圆弧对对称轴的转动惯量）
相互呼应，共同刻画了均匀圆弧的力学性质。

\end{document}
